\begin{abstract}
Emerging work like BitNet~\cite{bitnet} is ushering in a new generation of 1-bit Large Language Models (LLMs). This paper presents \textbf{\text{BitNet b1.58}}, a 1-bit LLM version in which all model parameters are restricted to the ternary set \{-1, 0, 1\}. This approach achieves parity with standard full-precision Transformer LLMs—using either FP16 or BF16 weights—in terms of perplexity and downstream task outcomes, provided model scale and training corpus remain unchanged. At the same time, it provides significant improvements in efficiency across multiple dimensions, including latency, memory usage, throughput, and power requirements. Importantly, the 1.58-bit LLM establishes a novel scaling law and a new methodological framework for developing subsequent high-performing yet resource-efficient LLM generations. In addition, it offers a fresh computational paradigm and paves the way for dedicated hardware architectures designed to leverage the properties of 1-bit LLMs.
\end{abstract}

\section{The Era of 1-bit LLMs}

The landscape of AI has witnessed exponential advances in both the scale and performance of Large Language Models (LLMs) in recent years. While these models have achieved unprecedented success across various natural language processing benchmarks, the continual expansion of model size introduces significant hurdles related to practical deployment, particularly with respect to energy expenditure and overall environmental and financial sustainability.

To mitigate these obstacles, one promising direction is post-training quantization, which constructs low-bitwidth LLM variants for inference~\cite{smoothquant, gptq, quip, quip_sharp}. This strategy compresses model weights and activations to lower precisions, sharply reducing demands on memory and computation. Recent industry practices have shifted from 16-bit representations down to more compact forms, such as 4-bit models~\cite{gptq, awq}. Despite their popularity in real-world LLM applications, post-training quantization remains an imperfect solution.

Recent advancements in 1-bit model design, exemplified by BitNet~\cite{bitnet}, have indicated a compelling strategy for reducing the operational costs of large language models (LLMs) without sacrificing their effectiveness. Standard LLMs utilize 16-bit floating-point representations (such as FP16 or BF16), and the primary computational burden is attributed to matrix multiplications dominated by floating-point addition and multiplication. BitNet, by contrast, employs integer addition in its matrix computations, thereby significantly lowering the associated energy demands. Because energy usage constitutes a fundamental constraint on performance for many hardware accelerators, these efficiency gains may also lead directly to faster computational throughput.

Beyond computation, inference costs are also shaped by the need to transfer model parameters from DRAM to on-chip accelerator memory (typically SRAM), which can become a bottleneck. Although expanding SRAM capacity has been explored as a way to boost throughput, this approach tends to incur considerably higher expenses than simply using DRAM. In contrast, 1-bit LLMs dramatically decrease the model's memory demands—both in terms of storage size and bandwidth—thus minimizing the cost and time required to move weights from DRAM, which enhances inference speed and efficiency.

Within this context, we propose a notable 1-bit LLM extension called \textbf{\text{BitNet b1.58}}, which restricts each parameter to ternary values, specifically from the set \{-1, 0, 1\}. By incorporating 0 into the original 1-bit BitNet framework, the model now encodes parameters in 1.58 bits in binary representation. \text{BitNet b1.58} maintains all of the computational advantages inherent to its predecessor—most notably, the novel computation paradigm that nearly eliminates multiplication operations in matrix multiplication, facilitating extensive hardware optimization. Furthermore, its energy profile remains on par with the original 1-bit BitNet, and it provides pronounced gains over FP16 LLMs regarding memory consumption, throughput, and latency.

Additionally, \text{BitNet b1.58} introduces two new benefits. First, by enabling explicit feature filtering through the use of 0-valued weights, it offers enhanced modeling power, which yields marked improvements over previous 1-bit LLMs. Second, empirical results demonstrate that \text{BitNet b1.58} attains perplexity and downstream task performance comparable to FP16 models beginning at the 3B parameter scale, given equivalent model configurations (e.g., model size, number of training tokens).

\section{BitNet b1.58}

\text{BitNet b1.58} is constructed upon the BitNet framework, utilizing a Transformer design that substitutes \emph{nn.Linear} layers with \emph{BitLinear} ones. Training is conducted from the ground up, employing 1.58-bit quantization for weights and 8-bit quantization for activations. Several updates distinguish this version from the initial BitNet, as outlined below.

\paragraph{Quantization Function.} To restrict weight values to the discrete set \{-1, 0, +1\}, we employ the \emph{absmean} quantization method. This approach begins by dividing the weight matrix by its mean absolute value, after which each element is rounded to the nearest integer within \{-1, 0, +1\}:

\begin{equation}
    \widetilde{W} = \mathrm{RoundClip}(\frac{W}{\gamma + \epsilon}, -1, 1),
\end{equation}

\begin{equation}
    \mathrm{RoundClip}(x, a, b) = \max(a, \min(b, \mathrm{round}(x))),
\end{equation}

\begin{equation}
    \gamma = \frac{1}{nm}\sum_{ij} |W_{ij}|.
\end{equation}

For activations, the quantization function mirrors the procedure used in BitNet. However, unlike the original, we omit scaling activations to the interval $[0, Q_b]$ prior to non-linearity. Instead, each token’s activation is scaled to $[-Q_b, Q_b]$, thereby eliminating the need for zero-point quantization. This modification simplifies both the implementation and system-level optimization, with our experiments indicating minimal impact on overall performance.

\paragraph{LLaMA-alike Components.}

Given the widespread adoption of LLaMA~\cite{llama, llama2} as a backbone in open-source LLM development, \text{BitNet b1.58} integrates similar architectural elements. The model leverages RMSNorm~\cite{rmsnorm}, SwiGLU~\cite{swiglu}, rotary positional encodings~\cite{rope}, and excludes biases throughout the network. This ensures seamless compatibility with leading open-source ecosystems such as Huggingface, vLLM~\cite{vllm}, and llama.cpp\footnote{\url{https://github.com/ggerganov/llama.cpp}}, requiring minimal adjustments for integration.

\section{Results}

We conducted a comparison between \text{BitNet b1.58} and our re-implemented FP16 \text{LLaMA LLM} across different model scales. Both models were pre-trained using the RedPajama dataset~\cite{redpajama} for a total of 100 billion tokens to guarantee a controlled evaluation environment. To assess generalization capabilities, we measured zero-shot accuracy over a suite of established language understanding tasks: ARC-Easy~\cite{arc}, ARC-Challenge~\cite{arc}, Hellaswag~\cite{hellaswag}, Winogrande~\cite{winoGrande}, PIQA~\cite{piqa}, OpenbookQA~\cite{openbookqa}, and BoolQ~\cite{boolq}. Additionally, we present validation perplexity on the WikiText2~\cite{wikitext2} and C4~\cite{c4} corpora.

For efficiency analysis, we benchmarked both GPU runtime memory usage and inference latency of \text{LLaMA LLM} and \text{BitNet b1.58}. These measurements were collected using the FasterTransformer\footnote{\url{https://github.com/NVIDIA/FasterTransformer}} library, optimized for rapid LLM inference on GPUs. For \text{BitNet b1.58}, the 2-bit kernel provided by Ladder~\cite{ladder} was incorporated. Our reporting focuses on the time to generate each output token, the primary bottleneck in inference operations.

The outcomes for perplexity and computational costs are compiled in Table~\ref{tab:ppl}. Notably, \text{BitNet b1.58} achieves comparable perplexity to the full precision \text{LLaMA LLM} when scaled to a 3B parameter model, while offering a 2.71$\times$ speedup and reducing GPU memory demands by 3.55$\times$. Particularly, with a 3.9B parameter count, \text{BitNet b1.58} demonstrates 2.4$\times$ faster inference, requires 3.32$\times$ less memory, and substantially surpasses the performance of the \text{LLaMA LLM} 3B model.

Detailed zero-shot performance on downstream tasks is provided in Table~\ref{tab:zero-shot}. The evaluation utilized the \emph{lm-evaluation-harness}\footnote{\url{https://github.com/EleutherAI/lm-evaluation-harness}} suite. Analysis indicates the performance disparity between \text{BitNet b1.58} and \text{LLaMA LLM} diminishes as model size increases, with \text{BitNet b1.58} matching full-precision baselines from the 3B scale upward. Consistent with perplexity findings, downstream task results highlight that \text{BitNet b1.58} 3.9B surpasses the 3B \text{LLaMA LLM} in both accuracy and efficiency. These results illustrate that \text{BitNet b1.58} achieves a Pareto advantage relative to leading LLMs.

\paragraph{Memory and Latency}

We extended our experiments to larger model variants, scaling up to 7B, 13B, and 70B parameters, and analyzed the associated computational costs. As depicted in Figure~\ref{fig:latency-memory-energy}, both latency and memory requirements display scaling trends, with acceleration becoming more pronounced as model size increases. Notably, at 70B parameters, \text{BitNet b1.58} delivers a speed-up of 4.1$\times$ relative to the \text{LLaMA LLM} baseline. This improvement primarily stems from the increasing proportion of time dedicated to \emph{nn.Linear} operations as models grow. Memory usage follows an analogous pattern since the full-precision embedding layer constitutes a smaller memory footprint at larger scales. Both latency and memory measurements utilize a 2-bit kernel implementation, suggesting that further reductions in cost may be achievable with additional optimization.

\paragraph{Energy}

We also estimated the energy expenditure associated with arithmetic operations in both \text{BitNet b1.58} and \text{LLaMA LLM}, emphasizing the matrix multiplication stage, which dominates overall LLM computational cost. Figure~\ref{fig:energy} breaks down the energy consumption profile for each approach. For \text{BitNet b1.58}, most energy is devoted to INT8 addition, whereas \text{LLaMA LLM} expends energy on both FP16 additions and multiplications. Drawing upon the energy analysis model in~\cite{energycost, pokebnn}, \text{BitNet b1.58} achieves a 71.4-fold reduction in energy consumption from arithmetic operations for matrix multiplication on 7nm hardware. End-to-end energy measurements with sequences of 512 tokens further indicate that scaling model size enhances the energy efficiency advantage of \text{BitNet b1.58} relative to the FP16 \text{LLaMA LLM} baseline. This arises because the contribution of \emph{nn.Linear} operations increases with model scale, while other sources of computational cost become less significant in larger architectures.

\paragraph{Throughput}

We evaluate the throughput of \text{BitNet b1.58} against the \text{LLaMA LLM} with 70B parameters, utilizing two 80GB A100 GPUs. Pipeline parallelism~\cite{gpipe} is employed to enable execution of the \text{LLaMA LLM} 70B model. Batch sizes are increased progressively until GPU memory is fully utilized, while maintaining a sequence length of 512. As summarized in Table~\ref{tab:throughput}, the \text{BitNet b1.58} 70B model is capable of handling batch sizes up to 11 times greater than those of \text{LLaMA LLM}, leading to an 8.9-fold improvement in throughput.

\textbf{\text{BitNet b1.58} introduces a novel scaling law relating model effectiveness to inference expenditure}. For comparison, the following correspondences between 1.58-bit and 16-bit models can be drawn based on the outcomes presented in Figure~\ref{fig:latency-memory-energy} and \ref{fig:energy}:

\begin{itemize}
    \item A 13B BitNet b1.58 model is more efficient across latency, memory consumption, and energy usage than a 3B FP16 LLM.
    \item A 30B BitNet b1.58 model demonstrates greater efficiency in latency, memory footprint, and energy requirements than a 7B FP16 LLM.
    \item A 70B BitNet b1.58 model surpasses a 13B FP16 LLM in terms of latency, memory, and energy efficiency.
\end{itemize}

\paragraph{Training with 2T Tokens}

The volume of training tokens is vital to the success of LLMs. To assess the token scalability of \text{BitNet b1.58}, we trained it on 2T tokens according to the data protocol of StableLM-3B~\cite{StableLM-3B-4E1T}, which represents the leading open-source model at the 3B scale. Both models underwent evaluation on a composite benchmark including Winogrande~\cite{winoGrande}, PIQA~\cite{piqa}, SciQ~\cite{sciq}, LAMBADA~\cite{lambada}, and ARC-easy~\cite{arc}. We present the zero-shot performance metrics in Table~\ref{tab:2t}. For benchmarks assessed with both accuracy and normalized accuracy, their mean is reported. Results for StableLM 3b at 2T tokens are cited directly from its technical documentation. Our analysis demonstrates that \text{BitNet b1.58} achieves superior outcomes on all evaluated tasks, highlighting that 1.58-bit LLMs are also capable of strong generalization.

\section{Discussion and Future Work}

\textbf{1-bit Mixture-of-Experts (MoE) LLMs}

Mixture-of-Experts (MoE) architectures have emerged as an efficient solution for scaling large language models (LLMs) in a cost-conscious manner. Although MoE substantially lessens the computational demand in terms of FLOPs, its significant memory usage and the resulting overhead in inter-device communication often impede practical deployment. These obstacles can be alleviated by employing 1.58-bit LLMs. Primarily, the lower memory requirements of such models reduce the number of hardware devices necessary to run MoE configurations. Furthermore, transferring activations across hardware nodes becomes much less burdensome. In an optimal scenario, if all model components fit within a single device, inter-chip communication could be eliminated entirely.

\textbf{Native Support of Long Sequence in LLMs}

The capability to process extended sequences is increasingly vital in the development of LLMs. A major limitation when scaling sequence length for inference is the substantial memory consumption associated with storing KV caches. The introduction of \text{BitNet b1.58} marks substantial progress in this area: by compressing activation storage from 16-bit to 8-bit representation, the maximum manageable context length is effectively doubled for fixed memory resources. Looking forward, additional reductions to 4 bits or potentially even less for 1.58-bit LLMs remain promising avenues for further research.

\textbf{LLMs on Edge and Mobile}

Deploying 1.58-bit LLMs could notably elevate the performance of language models operating on edge and mobile hardware. Devices in these categories frequently confront constraints in memory capacity and computational throughput, which limit both LLM size and efficiency. However, the reduced resource demands of 1.58-bit LLMs pave the way for deploying more capable models within these limited environments, making previously unattainable applications feasible. As a result, the computational utility and versatility of edge and mobile systems may be substantially enhanced. Additionally, since 1.58-bit LLMs are inherently more compatible with CPUs—the predominant processors in such platforms—models such as \text{BitNet b1.58} can execute efficiently and at scale, further extending their practical value in mobile and edge computing contexts.

\textbf{New Hardware for 1-bit LLMs}

Innovations such as those exemplified by Groq\footnote{\url{https://groq.com/}} highlight the prospects for custom-designed hardware, like language processing units (LPUs), tailored to LLM workloads. Building on this momentum, we advocate for the development of new hardware systems specifically designed to maximize the benefits of 1-bit LLMs, taking advantage of the distinct computation paradigm introduced by BitNet~\cite{bitnet}.

\end{document}